\section[Monge-Amp{\`e}re systems and Lagrangian pairs]{Monge--Amp{\`e}re systems and Lagrangian pairs}
\label{Monge-Ampere systems and Lagrangian pairs}

We start with the general def\/inition of Monge--Amp{\`e}re systems
\cite{B-G-G,L, Mo1,Mo2}.
Recall that a~contact structure $D$ on a~mani\-fold~$M$
is a subbundle of~$TM$ of codimension one locally def\/ined by a~contact
$1$-form $\theta$ by $D = \{ \theta = 0\}$ such that~$d\theta$ is non-degenerate,
that is, symplectic, on the bundle~$D$.
A~manifold endowed with a contact structure is called a contact manifold.
It is known that the dimension of a contact manifold is odd.

Let $(M, D)$ be a contact manifold of dimension $2n+1$ with
the contact structure \mbox{$D \subset TM$}.
A~{\it Monge--Amp{\`e}re system}
on $M$ is by def\/inition an exterior dif\/ferential system ${\mathcal M} \subset \Omega_M$
ge\-ne\-rated locally by a contact form $\theta$ for $D$ and an
$n$-form $\omega$ on~$M$: for each point $x \in M$,
there exists an open neighborhood~$U$ of~$x$ in~$M$ such that,
algebraically,
\begin{gather*}
{\mathcal M}\vert_U =
\langle \theta, d\theta, \omega
\rangle_{\Omega_U}.
\end{gather*}
Here $\Omega_M$ (resp.\ $\Omega_U$)
is the sheaf of germs of
exterior dif\/ferential forms
on $M$ (resp.\ on~$U$). In this case we call
$\omega$ a {\it local generator} of the Monge--Amp{\`e}re system
${\mathcal M}$ (modulo the contact ideal $\langle \theta, d\theta \rangle$).
Note that one may assume that $\omega$ is ef\/fective, i.e.,
$d\theta \wedge \omega \equiv
0\ (\rm{mod}\ \theta)$.
Also note that the $(n+1)$-form
$d\omega$ belongs to the contact ideal locally necessarily (see \cite{B1,B2, L}).

Let $D \subset TM$ be a vector bundle of rank $2n$ in the tangent bundle of a manifold $M$.
Recall that a conformal symplectic structure on $D$ is a reduction of the structure group of $D$
to the conformal symplectic group ${\rm CSp}({\mathbf{R}}^{2n})$.
If $(M, D)$ is a contact manifold of dimension $2n+1$, then
the conformal symplectic structure on $D$ is def\/ined locally by $d\theta$ for a contact form
$\theta$ which gives $D$ locally.
In particular, for each point $x \in M$, $D_x$ has the symplectic structure which is determined uniquely up to
a multiplication of a non-zero constant.
We call a linear subspace $W \subset D_x$
{\it Lagrangian} if $W$ is isotropic for the conformal symplectic structure on $D_x$ and $\dim_{{\mathbf{R}}} W = n$.
A subbundle $E \subset D$ is called a {\it Lagrangian subbundle} if $E_x \subset D_x$ is Lagrangian for any $x \in M$.

Now we def\/ine the key notion in this paper.

\begin{Def}
Let $(M, D)$ be a contact manifold.
A {\it Lagrangian pair} is a pair $(E_1, E_2)$
of Lagrangian subbundles of $D$ with respect to the conformal symplectic structure
on $D$ which satisf\/ies the condition $D = E_1 \oplus E_2$.
\end{Def}

\begin{Rem}
In \cite[\S~5.2]{B}, the notion of bi-Lagrangian structure is def\/ined as the transverse pair of Lagrangian foliations in a symplectic manifold.
Since we treat the contact case, it might be natural to use
the terminology  ``Legendrian'' instead of ``Lagrangian''.
However we would like to use ``Lagrangian'' in the general cases, and
to use the terminology ``Legendrian'' just for the integrable cases, such as ``Legendrian submanifolds'' and
``Legendrian f\/ibrations'' (see Section~\ref{Hesse representations.}).
\end{Rem}

The standard example of Lagrangian pair is given as follows.

{\bf The standard example.}
The standard example of Lagrangian pair is given on
the standard Darboux model.
Consider ${\mathbf{R}}^{2n+1}$ with coordinates
$(x_1,\dots,x_n,z,p_1,\dots,p_n)$ and with the standard
contact structure
$D_{\rm{st}} = \{ v \in TM \,|\, \theta_{\rm{st}}(v) = 0\}$ def\/ined by
the standard contact form
$\theta_{\rm{st}} = dz - \sum\limits_{i=1}^np_idx_i$.
Then we set
\begin{gather*}
E^{\rm{st}}_1   =
\{ v \in D_{\rm{st}} \,|\,
dp_1(v) = \cdots = dp_n(v) = 0 \} = \left\langle \frac{\partial}{\partial x_1} +
p_1\frac{\partial}{\partial z},
\dots, \frac{\partial}{\partial x_n} + p_n\frac{\partial}{\partial z}\right\rangle,
\\
E^{\rm{st}}_2  =
\{ u \in D_{\rm{st}} \,|\,
dx_1(u) = \cdots = dx_n(u) = 0\} =
\left\langle \frac{\partial}{\partial p_1},
\dots, \frac{\partial}{\partial p_n}\right\rangle.
\end{gather*}
Then $(E^{\rm{st}}_1, E^{\rm{st}}_2)$ is a Lagrangian pair on
$({\mathbf{R}}^{2n+1}, D_{\rm{st}})$.


Note that,
in \cite{Tak}, Takeuchi called
a contact structure $D$ endowed with a Lagrangian pair $(E_1,E_2)$
a {\it Lagrangian contact structure} $(D; E_1,E_2)$ and
gave a detailed study on this geometric structure.

Moreover we consider an exterior dif\/ferential system associated to a given
Lagrangian pair or Lagrangian contact structure.

\begin{Def}
A Monge--Amp{\`e}re system $\mathcal{M}$ is called
a {\it decomposable Monge--Amp{\`e}re system with a Lagrangian pair} $(E_1,E_2)$
if in a neighborhood of each point of $M$, there exists a local generator $\omega$
of $\mathcal{M}$
satisfying the following {\it decomposing condition}:
\begin{gather*}
i_u\omega = 0\ (u\in E_1), \quad \text{and} \quad
\omega\vert_{E_2}\ \text{is a volume form on}\ E_2.
\end{gather*}
\end{Def}

Such a decomposable Monge--Amp{\`e}re system
of Lagrangian type is a decomposable Monge--Amp{\`e}re system with
the characteristic system $E_1$ in the sense of~\cite{M-M}.
Note that this class of Monge--Amp{\`e}re equations have been introduced in~\cite{G}.
Also note that decomposable Monge--Amp{\`e}re systems are studied from another perspective in~\cite{A-B-M-P} by the name of Goursat equations.


\begin{Def}
A Monge--Amp{\`e}re system $\mathcal{M}$ is called
a {\it bi-decomposable Monge--Amp{\`e}re system with a Lagrangian pair} $(E_1,E_2)$
if, around each point of $M$, there exists a local generator~$\omega$
of~$\mathcal{M}$
of the form
\begin{gather*}
\omega=\omega_1-\omega_2
\end{gather*}
by $n$-forms $\omega_1$ and $\omega_2$
satisfying the following {\it bi-decomposing condition}:
\begin{gather*}
i_u\omega_1=0\quad (u\in E_2),\qquad i_v\omega_2=0\quad (v\in E_1),
\\
\omega_1| _{E_1}\ \text{is a volume form on}\ E_1,\quad
\mathrm{and}\quad \omega_2| _{E_2}\ \text{is a volume form on}\ E_2.
\end{gather*}
We call such an $n$-form $\omega$ a {\it bi-decomposable form} and
$(\omega_1, \omega_2)$ a {\it bi-decomposition} of~$\omega$.
Given a~Lagrangian pair $(E_1,E_2)$ on $(M, D)$ and $n$-forms $\omega_1$,
$\omega_2$ satisfying the bi-decomposing condition for $(E_1,E_2)$,
we def\/ine a Monge--Amp{\`e}re system
${\mathcal M}$ with Lagrangian pair
by setting $\omega = \omega_1-\omega_2$.
\end{Def}


\begin{Rem}
An immersion $f\colon  L \to M$ of an $n$-dimensional manifold $L$ to $M$
is called a~{\it geometric solution} of a Monge--Amp{\`e}re system
${\mathcal M} = \langle \theta, d\theta, \omega\rangle$ if
$f^*{\mathcal M} = 0$, namely, if $f^*\theta = 0$, i.e., $f$ is a Legendrian
immersion, and $f^*\omega = 0$.

Any geometric solution to
a decomposable (resp.\ bi-decomposable) Monge--Amp{\`e}re system~${\mathcal M}$
with a Lagrangian pair $D = E_1\oplus E_2$
on a contact manifold $(M^{2n+1}, D)$ has a crucial property.

In fact, an immersion $f\colon L \to M$ is a geometric solution of a decomposable Monge--Amp{\`e}re system
with a~Lagrangian pair $D = E_1\oplus E_2$ if and only if $f_*(T_pL) \cap (E_1)_{p} \not= \{ 0\}$.

For a bi-decomposable Monge--Amp{\`e}re system ${\mathcal M}$
with a Lagrangian pair $D = E_1\oplus E_2$,
suppose that $f\colon L \to M$ is a geometric solution of ${\mathcal M}$.
Then, for any $p \in L$, we see that
\begin{gather*}
E_1 \cap f_*(T_pL) = \{ 0\},\qquad \text{if and only if}\quad E_2 \cap f_*(T_pL) = \{ 0\},
\end{gather*}
equivalently,
\begin{gather*}
E_1 \cap f_*(T_pL) \not= \{ 0\},\qquad \text{if and only if}\quad E_2 \cap f_*(T_pL) \not= \{ 0\}.
\end{gather*}
In fact, let $W$ be an $n$-plane
in $D_p = (E_1)_p \oplus (E_2)_p$ satisfying
$\omega\vert_W = 0$.
The direct sum decomposition
def\/ines projections $\pi_1\colon D_p \to (E_1)_p$ and $\pi_2\colon D_p \to (E_2)_p$.
Then $E_1 \cap W = \{ 0\}$ if and only if $\pi_2\vert_{W}\colon W \to
(E_2)_p$ is an isomorphism, and
$E_2 \cap W = \{ 0\}$ if and only if $\pi_1\vert_{W}\colon W \to
(E_1)_p$ is an isomorphism, respectively.
For any bi-decomposition $\omega = \omega_1 - \omega_2$ and
for any basis $u_1, \dots, u_n$ of $W$, we have
\begin{gather*}
\omega_1(\pi_1(u_1), \dots, \pi_1(u_n)) =
\omega_1(u_1, \dots, u_n) =
\omega_2(u_1, \dots, u_n) =
\omega_2(\pi_2(u_1), \dots, \pi_2(u_n)).
\end{gather*}
Using the bi-decomposing condition again,
we see that $\pi_1\vert_{W}$ is an isomorphism if and only if
the most left hand side is non-zero, and it is equivalent to
the condition that
$\pi_2\vert_{W}$ is an isomorphism.
\end{Rem}



Now we are led to natural questions:
\begin{itemize}\itemsep=0pt
\item Is the Lagrangian pair $(E_1,E_2)$ uniquely determined
by a decomposable form?

\item Is the Lagrangian pair $(E_1,E_2)$ uniquely determined
by a bi-decomposable form?

\item Is the Lagrangian pair $(E_1, E_2)$ recovered
only from the Monge--Amp{\`e}re system $\mathcal{M}$?
\end{itemize}

As is stated in Introduction, the f\/irst question is answered negatively.
To answer the second and third questions, we recall the basic def\/initions.

\begin{Def}
\label{isomorphism}
Let $(M, D)$,  $(M', D')$ be contact manifolds of dimension $2n+1$,
and ${\mathcal M}$, ${\mathcal M'}$ Monge--Amp{\`e}re systems
on contact manifolds $(M,D)$, $(M',D')$ respectively.
A dif\/feomorphism $\Phi\colon M \longrightarrow M'$ is called
an {\it isomorphism of Monge--Amp{\`e}re systems}
if (1)~$\Phi$ is a~contactomorphism,
namely $(\Phi_*)D=D'$, and
(2)~${\Phi}^*{\mathcal M'}={\mathcal M}$.
\end{Def}

Now suppose that contact manifolds $(M, D)$,  $(M', D')$ of dimension
$2n+1$ are endowed with
Lagrangian pairs $(E_1, E_2)$, $(E'_1, E'_2)$ respectively,
namely, that the decompositions $D = E_1\oplus E_2$ and $D' = E'_1\oplus E'_2$
are given. Suppose $n \geq 3$.
Then, from the result in
Section~\ref{Lagrangian pair and bi-decomposable forms.}
mentioned above, we have that
any isomorphism ${\Phi}$ of a Monge--Amp{\`e}re system
${\mathcal M}$ with the Lagrangian pair $(E_1, E_2)$ and
a Monge--Amp{\`e}re system
${\mathcal M}'$ with the Lagrangian pair $(E'_1, E'_2)$
necessarily preserves the Lagrangian pairs up to ordering, namely,
$(\Phi_*)E_1 = E'_1$, $(\Phi_*)E_2 = E'_2$  or
$(\Phi_*)E_1 = E'_2$, $(\Phi_*)E_2 = E'_1$.

In the case $n = 2$, a Lagrangian pair is not uniquely determined
from a Monge--Amp{\`e}re system and moreover, the automorphism pseudo-group
may be of inf\/inite dimension. For instance, consider the equation
${\rm Hess} = -1$ which is isomorphic to the linear wave
equation $f_{x_1x_1} - f_{x_2x_2} = 0$ and to the equation
$f_{x_1x_2} = 0$. Then the last equation has inf\/inite-dimensional
automorphisms induced by dif\/feomorphisms $(x_1, x_2)
\mapsto (X_1(x_1), X_2(x_2))$.

\section{Lagrangian pair and bi-decomposable form}
\label{Lagrangian pair and bi-decomposable forms.}

Let $(M, D)$ be a contact manifold of dimension $2n+1$ with
a contact structure $D \subset TM$.
We have def\/ined in Section~\ref{Monge-Ampere systems and Lagrangian pairs} the notion of bi-decomposing
conditions and bi-decomposable forms on~$(M, D)$.
Then, f\/irst we show

\begin{Lem}
\label{bi-decomposability}
Let $\omega$ be an $n$-form on a contact manifold $(M, D)$,
$D = \{ u \in TM \,|\, \theta(u) = 0\}$
for a local contact form $\theta$ defining $D$,
and $(E_1, E_2)$
a Lagrangian pair of~$D$.
Assume that $\omega$ is
a~bi-decomposable form for $(E_1, E_2)$,
and $\omega = \omega_1 - \omega_2$ is
any bi-decomposition of $\omega$ for $(E_1, E_2)$.
Then locally there exists a coframe
$\theta, \alpha_1, \dots, \alpha_n, \beta_1, \dots, \beta_n$
of $T^*M$ such that
\begin{gather*}
E_1 = \{ v \in D \,|\, \beta_1(v) = \cdots = \beta_n(v) = 0\}, \qquad
E_2 = \{ u \in D \,|\, \alpha_1(u) = \cdots = \alpha_n(u) = 0\},
\end{gather*}
and that the $n$-forms
\begin{gather*}
\widetilde{\omega}_1 = \alpha_1\wedge \cdots \wedge \alpha_n,
\qquad
\widetilde{\omega}_2 = \beta_1\wedge \cdots \wedge \beta_n
\end{gather*}
satisfy the bi-decomposing condition for
$(E_1, E_2)$ with
\begin{gather*}
\widetilde{\omega}_1 \equiv \omega_1, \qquad
\widetilde{\omega}_2 \equiv \omega_2, \qquad
\omega \equiv \widetilde{\omega}_1 - \widetilde{\omega}_2
\end{gather*}
up to a multiple of $\theta$.
\end{Lem}

\begin{proof}
The proof is based on the fact that the symplectic group on a f\/inite-dimensional symplectic vector space
acts transitively on the set of transversal pairs of Lagrangian subspaces.

Let $X_1, \dots, X_n$ and $P_1, \dots, P_n$ be local frames
of~$E_1$ and~$E_2$ respectively.
Let~$R$ be the Reeb vector f\/ield
for a local contact form $\theta$ def\/ining $D$. Recall that~$R$ is def\/ined uniquely by $i_R \theta = 1$, $i_R d\theta = 0$.
Consider
the dual coframe $\theta, \alpha_1,\dots,\alpha_n, \beta_1,\dots,\beta_n$
of~$T^*M$
to the frame~$R$, $X_1, \dots, X_n$, $P_1, \dots, P_n$ of~$TM$.
Then we see, from the bi-decomposing condition,
that there exist an $(n-1)$-form $\gamma$ and a~non-vanishing function $\mu$ on $M$ such that
$\omega_1 = \mu(\alpha_1 \wedge \cdots \wedge \alpha_n) + \theta\wedge\gamma$.
By replacing $\alpha_1$ by $\frac1{\mu}\alpha_1$,
we may suppose $\mu \equiv 1$.
Similarly, we have $\omega_2 \equiv
\beta_1 \wedge \cdots \wedge \beta_n \ \rm{mod}\ \theta$.
\end{proof}


Next we show that $\mathcal{M}$ has the unique
local bi-decomposable generator $\omega$ modulo $\theta$.

\begin{Th}
\label{bi-decomposable-form}
Let $(E_1, E_2)$ be a Lagrangian pair and
$\omega$, $\omega'$ be two bi-decomposable $n$-forms
for the Lagrangian pair $(E_1, E_2)$ on $(M, D)$.
Assume that they generate the same Monge--Amp{\`e}re system
\begin{gather*}
{\mathcal M} =
\langle \theta, d\theta, \omega \rangle
=
\langle \theta, d\theta, \omega' \rangle.
\end{gather*}
Then there exist
locally a non-vanishing function $\mu$ and an $(n-1)$-form $\eta$
on $M$ such that $\omega' = \mu \omega + \theta \wedge \eta$.
\end{Th}

To show Theorem~\ref{bi-decomposable-form},
we study, for each $x \in M$, the symplectic exterior linear algebra
on the symplectic vector space $V = D_x$ with the symplectic form
$\Theta = d\theta\vert_{D_x}$ and with the decomposition
$V = V_1\oplus V_2$, $V_1 = (E_1)_x$, $V_2 = (E_2)_x$,
of $(V, \Theta)$ into Lagrangian subspaces.


Let $(V,\Theta)$ be a $2n$-dimensional symplectic vector space.
We say that an $n$-form $\omega \in \wedge^n V^*$ is
{\it bi-decomposable} if
there exist a decomposition $V = V_1\oplus V_2$ of~$V$ into
Lagrangian sub\-spa\-ces~$V_1$,~$V_2$ in~$V$ and
$n$-forms $\omega_1, \omega_2 \in \wedge^n V^*$ such
that $\omega = \omega_1 - \omega_2$,
$i_u\omega_1 = 0$ $(u \in V_2)$, $i_v\omega_2 = 0$ $(v \in V_1)$,
$\omega_1\vert_{V_1} \not= 0$, $\omega_2\vert_{V_2} \not= 0$.
In this case $(\omega_1, \omega_2)$ is called a {\it bi-decomposition}
of~$\omega$.
Then, similarly as the proof of Lemma~\ref{bi-decomposability}, we have that
there exists a basis
$\{ a_1,\dots, a_n, b_1,\dots, b_n\}$ of~$V$
such that
\begin{gather*}
\omega_1 = a_1^*\wedge \cdots \wedge a_n^*, \qquad
\omega_2 = b_1^*\wedge \cdots \wedge b_n^*,
\end{gather*}
and that
\begin{gather*}
V_1= \langle a_1,\dots, a_n\rangle,
\qquad
V_2= \langle b_1,\dots, b_n\rangle,
\end{gather*}
where $\{ a_1^*,\dots, a_n^*, b_1^*,\dots, b_n^*\}$ denotes the dual basis
of $\{ a_1,\dots, a_n, b_1,\dots, b_n\}$.
Note that
$V_2$ (resp.~$V_1$) coincides with the annihilator of $a_1^*,\dots, a_n^*$
(resp.~$b_1^*,\dots, b_n^*$).
Moreover we see that
there exist a symplectic basis $\{ a_1,\dots, a_n, b_1,\dots, b_n\}$
of $(V, \Theta)$
and a non-zero constants~$a$,~$b$ such that
\begin{gather*}
\omega =
a   a_1^*\wedge \cdots \wedge a_n^* - b   b_1^*\wedge \cdots \wedge b_n^*.
\end{gather*}
If we replace $a_1$, $b_1$ by $(1/b) a_1$, $b b_1$
respectively,
so $a_1^*$, $b_1^*$ by $b a_1^*$, $(1/b) b_1^*$,
and set $c = ab$, then we have the following:

\begin{Lem}
\label{symplectic basis}
Let $\omega \in \wedge^n V^*$ be a bi-decomposable $n$-form on
a $2n$-dimensional symplectic vector space $(V, \Theta)$, and
$\omega = \omega_1 - \omega_2$ a~bi-decomposition of~$\omega$.
Then there exist a symplectic basis
$\{ a_1,\dots, a_n, b_1,\dots, b_n\}$ of~$(V, \Theta)$
and a non-zero constant $c$ such that
\begin{gather*}
\omega_1  =
c  a_1^*\wedge \cdots \wedge a_n^*, \qquad
\omega_2  =   b_1^*\wedge \cdots \wedge b_n^*,
\\
\omega  =  c  a_1^*\wedge \cdots \wedge a_n^* -
b_1^*\wedge \cdots \wedge b_n^*,
\qquad
\Theta  =  a_1^* \wedge b_1^* + \cdots + a_n^* \wedge b_n^*.
\end{gather*}
\end{Lem}

A form $\varphi \in \wedge^n V^*$ on
a symplectic vector space $(V,\Theta)$
is called {\it effective} if
the interior product $i_{X_{\Theta}}\varphi = 0$ for
the $2$-vector $i_{X_{\Theta}}$ corresponding to the
$2$-form~$\Theta$.
Note that
\begin{gather*}
X_{\Theta} = a_1 \wedge b_1 + \cdots + a_n \wedge b_n
\end{gather*}
in terms of the basis in Lemma \ref{symplectic basis}.
That $\varphi \in \wedge^n V^*$ is ef\/fective if and only if
the wedge $\varphi\wedge\Theta$ with the symplectic form~$\Theta$ is equal to zero~\cite{B2, L}.
Then we have, by Lemma~\ref{symplectic basis}:

\begin{Cor}
\label{effective}
If $\omega \in \wedge^n V^*$ is bi-decomposable and
$\omega = \omega_1 - \omega_2$ is any bi-decomposition, then~$\omega$,~$\omega_1$ and~$\omega_2$ are all effective.
\end{Cor}

We will apply the following basic result to our situation.

\begin{Lem}[\protect{\cite[Theorem~1.6]{L}}]\label{Lych}
Let $\omega$, $\omega'$ be effective $k$-forms on a symplectic vector space
$(V, \Theta)$, $(0\leq k\leq n)$.
Suppose that, for every isotropic subspace $L\subset V$ on which
$\omega\vert_L = 0$, the form $\omega'$ also vanishes on~$L$,
$\omega'\vert_L = 0$.
Then we have $\omega' =\mu{\omega}$ for some~$\mu \in {\mathbf{R}}$.
\end{Lem}

Then we have the following:

\begin{Lem}
\label{bi-decomposable2}
Let $\omega$, ${\omega}'$ be bi-decomposable forms on $(V, \Theta)$.
Suppose that~${\omega}'$ is of the form
$\lambda \omega+\phi \wedge \Theta$
for a scalar $\lambda$ and an $(n-2)$-form $\phi$.
Then ${\omega}' = \mu\omega$ for a scalar~$\mu$.
\end{Lem}

\begin{proof}
By Corollary \ref{effective}, the $n$-form $\omega$ is ef\/fective
and so is the $n$-form~${\omega}'$.
For every Lagrangian subspace~$L$ $(\Theta\vert_L = 0)$ on which
$\omega\vert_L = 0$,
the form $\omega'$ also vanishes.
Therefore, by Lemma~\ref{Lych},
it follows that
$\omega' =\mu{\omega}$
for a scalar~$\mu$.
\end{proof}

\begin{Rem}
Lemma \ref{bi-decomposable2} can be shown, as an alternative proof, by
applying Lefschetz isomorphism $\Theta^2\colon
\wedge^{n-2} V \to \wedge^{n+2} V$. In fact, from ${\omega}' - \mu\omega = \phi \wedge \Theta$
we have $\phi\wedge \Theta^2 = 0$, which implies~$\phi = 0$.
\end{Rem}


\begin{proof}[Proof of Theorem \ref{bi-decomposable-form}]
Since $\omega'$ belongs to ${\mathcal M} = \langle \theta, d\theta, \omega
\rangle$,
we set $\omega' = \lambda\omega +  d\theta\wedge\phi + \theta\wedge \eta$,
for a~function~$\lambda$, an $(n-2)$-form $\phi$ and an $(n-1)$-form $\eta$
on~$M$.
For each $x \in M$, we have
$\omega'\vert_{D_x} = \lambda(x)\omega\vert_{D_x} +
\Theta\wedge\phi\vert_{D_x}$,
where $\Theta = d\theta\vert_{D_x}$.
Then, by Lemma~\ref{bi-decomposable2}, we have
$\omega'\vert_{D_x} = \mu(x)\omega\vert_{D_x}$ for a scalar~$\mu(x)$
depending on $x \in M$.
Since $\omega\vert_{(E_1)_x}$ is a volume form, we see that~$\mu(x)$
is unique and
of class $C^\infty$. Moreover there exists a $C^\infty$ $(n-1)$-form $\eta$
such that $\omega' - \mu\omega = \theta \wedge \eta$, which implies
the required consequence.
\end{proof}


Moreover we show the following basic result:
\begin{Th}\label{Uniqueness}
Assume that $n\geq 3$.
Let
$\mathcal{M}=
\langle \theta, d\theta, \omega \rangle$ be
a Monge--Amp{\`e}re system
with a~Lag\-rangian pair locally generated by a bi-decomposable
$n$-form~$\omega$. Then~$\mathcal{M}$
canonically determines the Lagrangian pair~$(E_1,E_2)$.
Namely, if $(\omega_1, \omega_2)$
$($resp.\ $(\omega'_1, \omega'_2))$
is a bi-decomposition of~$\omega$ for
a~Lag\-rangian pair $(E_1,E_2)$ $($resp.~$(E_1', E_2'))$
enjoying the bi-decomposing condition.
Then we have
\begin{gather*}
E_1'=E_1,\quad  E_2'=E_2, \qquad {\mathrm or} \qquad
E_1' = E_2, \quad  E_2' = E_1.
\end{gather*}
\end{Th}

Theorem \ref{Uniqueness} follows immediately from the following:

\begin{Prop}
\label{uniqueness 2}
Assume that $n\geq 3$.
Let $(V,\Theta)$ be a $2n$-dimensional symplectic vector space.
Let $\omega=\omega_1-\omega_2$ be a bi-decomposable
$n$-form for a Lagrangian pair $(V_1, V_2)$ of $V$.
Then the bi-decomposition $(\omega_1, \omega_2)$ of~$\omega$
is unique up to ordering:
If $\omega=\omega'_1-\omega'_2$ is another bi-decomposition
of $\omega$ for another Lagrangian pair $(V'_1, V'_2)$ of~$V$, then
\begin{gather*}
V'_1 = V_1, \qquad V'_2 = V_2, \qquad \omega'_1 = \omega_1, \qquad \omega'_2 = \omega_2,
\qquad {\mathrm{or}} \\
V'_1 = V_2, \qquad V'_2 = V_1, \qquad \omega'_1 = - \omega_2, \qquad \omega'_2 = - \omega_1.
\end{gather*}
\end{Prop}



The bi-decomposition of $\omega$ (Proposition~\ref{uniqueness 2})
is given by using the symplectic structure,
based on Hitchin's result \cite[Propositions 2.1,~2.2]{H},
in the case that $n$ is odd and $n \geq 3$, as follows:

Let $\omega = \omega_1 - \omega_2$, with
$\omega_1 = c a_1^*\wedge \cdots \wedge a_n^*$
and
$\omega_2 = b_1^*\wedge \cdots \wedge b_n^*$, as in
Lemma~\ref{symplectic basis}.
Let $\epsilon=a^*_1\wedge \cdots \wedge a^*_n
\wedge b^*_1\wedge \cdots b^*_n$
be the associated basis vector for $\Lambda^{2n}V^*$,
which is the
intrinsically def\/ined volume form by the symplectic structure on~$V$.
From the isomorphism $A\colon \Lambda^{2n-1}V^* \longrightarrow
V\otimes \wedge^{2n}V^*$ induced by the exterior product
$V^*\otimes\wedge^{2n-1}V^* \to \wedge^{2n}V^*$,
we def\/ine, for each $\psi \in \wedge^nV^*$,
a linear transformation
$K_{\psi}=K_{\psi}^{\epsilon}\colon V \longrightarrow V$ by
$K_{\psi}(u)\epsilon = A(i_u(\psi)\wedge \psi)$, and put
$\lambda(\psi)=\lambda_{\epsilon}(\psi)=
\frac{1}{2n}\operatorname{tr}(K_{\psi}^2)$.

In particular, we set $\psi = \omega$.
Then we have $K_{\omega}a_i=-c a_i$, $K_{\omega}b_i
= (-1)^{n+1} c b_i = cb_i$.
Then we have $K_{\omega}^2 = c^2\operatorname{id}$,
therefore $\lambda(\omega) = c^2$.
Since $K_{\omega}^*a_i^* = -ca_i^*$, $K_{\omega}^*b_i^* =cb_i^* $,
we have
\begin{gather*}
K_{\omega}^*\omega = c(-c)^n a^*_1\wedge \cdots \wedge a^*_n -
c^{n}b^*_1\wedge \cdots \wedge b^*_n =
(-c)^n(\omega_1 + \omega_2).
\end{gather*}
Thus, from
$\omega={\omega}_1-{\omega}_2$, $-\frac{1}{c^n}K_{\omega}^*\omega={\omega}_1+{\omega}_2$
and $\lambda(\omega) = c^2$ $(c >0, <0)$,
we see
\begin{gather*}
\omega_1 = \frac{1}{2}\big(\omega \mp \lambda(\omega)^{-\frac{n}{2}}
K_{\omega}^*\omega\big),\qquad
\omega_2= \frac{1}{2}\big({-}\omega \mp \lambda(\omega)^{-\frac{n}{2}}
K_{\omega}^*\omega\big).
\end{gather*}
Since $K_{\omega}^*\omega$ and $\lambda(\omega)$ are
intrinsically determined from the symplectic structure on~$V$,
so is the decomposition of~$\omega$.

To verify Proposition~\ref{uniqueness 2} in general case,
we observe a fact from projective geometry of Pl\"{u}cker embeddings.
Consider Grassmannian ${\rm Gr}(n, V^*)$ consisting of all
$n$-dimensional subspaces
in the $2n$-dimensional vector space $V^*$, and its
Pl\"{u}cker embedding ${\rm Gr}(n, V^*) \hookrightarrow P(\wedge^{n} V^*)$.

\begin{Lem}
\label{proj.geom}
Let $\Lambda_1, \Lambda_2 \in {\rm Gr}(n, V^*)$ with $\Lambda_1 \cap \Lambda_2 = \{ 0\}$.
\begin{enumerate}\itemsep=0pt
\item[$(1)$] The projective line in $P(\wedge^{n} V^*)$
through the two points~$\Lambda_1$,~$\Lambda_2$ does not intersect
with ${\rm Gr}(n, V^*)$ at other points.

\item[$(2)$] Assume $n \geq 3$. Let $\Lambda_3$ be another point in~${\rm Gr}(n, V^*)$
different from $\Lambda_1$, $\Lambda_2$. Then the projective plane
in $P(\wedge^{n} V^*)$ through the three points $\Lambda_1$, $\Lambda_2$,
$\Lambda_3$
does not intersect with ${\rm Gr}(n, V^*)$ at other points.
\end{enumerate}
\end{Lem}

\begin{proof}
Choose a basis $e_1, \dots, e_n$ of $\Lambda_1$ and
$e_{n+1}, \dots, e_{2n}$ of $\Lambda_2$ to form
a basis of~$V^*$. Let $( p_{i_1, i_2, \dots, i_n})$
denote Pl\"{u}cker coordinates of the Pl\"{u}cker embedding. Here
$i_1, i_2, \dots, i_n$ are $n$-tuple chosen from $\{ 1, 2, \dots, 2n\}$.
For any sequences $1 \leq j_1 < j_2 < \cdots < j_{n-1} \leq 2n$ of
$(n-1)$-letters and
$1 \leq k_1 < k_2 < \cdots < k_{n+1} \leq 2n$ of $(n+1)$-letters,
we have the Pl\"{u}cker relation
\begin{gather*}
\sum_{\ell = 1}^{n+1} (-1)^\ell   p_{j_1, j_2, \dots, j_{n-1}, k_{\ell}}
p_{k_1, k_2,
\dots, \breve{k}_\ell, \dots, k_{n+1}} = 0
\end{gather*}
(see \cite{G-K-Z,G-H} for instance).
To see (1),  take a point $W  \in {\rm Gr}(n, V^*)$
on the projective line through~$\Lambda_1$,~$\Lambda_2$.
Since Pl\"{u}cker coordinates of $\Lambda_1$ (resp.\ $\Lambda_2$)
are given by $p_{1,\dots,n} = 1$ (resp.\ $p_{n+1, \dots, 2n} = 1$)
with zero other coordinates, we have that Pl\"{u}cker coordinates of~$W$
are given by $p_{1,\dots,n} = \lambda$, $p_{n+1, \dots, 2n} = \mu$
for some~$\lambda$,~$\mu$, other coordinates $p_{i_1, \dots, i_n}$ being
null.
Then, by applying the Pl\"{u}cker relation for $(j_1, \dots, j_{n-1}) =
(1, 2, \dots, n-1)$ and $(k_1, k_2, \dots, k_{n+1}) = (n-1, n, \dots, 2n)$,
we have $\lambda\mu = 0$. Therefore $\lambda = 0$ or $\mu = 0$, namely,
$W = \Lambda_1$ or $W = \Lambda_2$.

To see (2), let $\bar{p}=(\bar{p}_{i_1, i_2, \dots, i_n})$
be coordinates of~$\Lambda_3$.
Then the coordinates $p=(p_{i_1, i_2, \dots, i_n})$
for a point~$W$ on the plane
through $\Lambda_1$, $\Lambda_2$, $\Lambda_3$ are given by
\begin{gather*}
p_{1,\dots,n} = \lambda + \bar{p}_{1,\dots,n}, \qquad
p_{n+1, \dots, 2n} = \mu + \bar{p}_{n+1, \dots, 2n}, \qquad
p_{i_1, \dots, i_n} = \bar{p}_{i_1, \dots, i_n},
\end{gather*}
$\{i_1, \dots, i_n\} \not= \{ 1, \dots, n\}, \{ n+1, \dots, 2n\}$.
Suppose $W \in {\rm Gr}(n, V^*)$, $W \not= \Lambda_1, \Lambda_2, \Lambda_3$.
Then we see that
both~$p$ and~$\bar{p}$ satisfy
the Pl\"{u}cker relations and that $\lambda\mu\not= 0$.
Let $\{ i_1, \dots, i_n \} \not= \{ 1, \dots, n\}, \{ n+1, \dots, 2n\}$.
Then, because $n \geq 3$,  there exist $\ell$, $\ell'$, $\ell \not= \ell'$,
such that $i_\ell \not\in \{ n+1, \dots, 2n\}$ and $i_{\ell'}
\not\in \{ n+1, \dots, 2n\}$ or that $i_\ell \not\in \{ 1, \dots, n\}$
and $i_{\ell'} \not\in \{ 1, \dots, n\}$. In the former case,
we take $(i_1, \dots, \breve{i}_k, \dots, i_n)$,  $(i_k, n+1, \dots, 2n)$,
and write the Pl\"{u}cker relation to get $\mu\cdot\bar{p}_{i_1, \dots, i_n}
= 0$.
In the latter case, we take $(i_1, \dots, \breve{i}_k, \dots, i_n),
(1, \dots, n, i_k)$ to get $\lambda\cdot\bar{p}_{i_1, \dots, i_n} = 0$.
Therefore we obtain that $\bar{p}_{i_1, \dots, i_n} = 0$
for $\{ i_1, \dots, i_n \} \not= \{ 1, \dots, n\}, \{ n+1, \dots, 2n\}$.
This means that $\Lambda_3$
lies on the projective line through~$\Lambda_1$,~$\Lambda_2$,
which contradicts~(1).
\end{proof}

\begin{proof}[Proof of Proposition \ref{uniqueness 2}]
Set $\omega_1 - \omega_2 = \omega'_1 - \omega'_2$. Then
$\omega_1 - \omega_2 - \omega'_1$ is equal to a decomposable form
$- \omega'_2$.
Let $\Lambda_1$, $\Lambda_2$, $\Lambda'_1$, $\Lambda'_2$ be
points in ${\rm Gr}(n, V^*) \subset P(\wedge^n V^*)$
corresponding to $\omega_1$, $\omega_2$, $\omega'_1$, $\omega'_2$ respectively.
Then the projective plane through $\Lambda_1$, $\Lambda_2$, $\Lambda'_1$
intersects with~${\rm Gr}(n, V^*)$ also at~$\Lambda'_2$.
Note that $\Lambda'_1 \not= \Lambda'_2$ as well as
$\Lambda_1 \not= \Lambda_2$.
Suppose $\Lambda'_1 \not= \Lambda_1, \Lambda_2$.
By Lemma~\ref{proj.geom}(2), we have
$\Lambda'_2 = \Lambda_1$ or $\Lambda'_2 = \Lambda_2$.
Then we have that~$\Lambda'_1$ lies on the line through $\Lambda_1$, $\Lambda_2$.
By Lemma~\ref{proj.geom}(1), we conclude that
$\Lambda'_1 = \Lambda_1$ or $\Lambda'_1 = \Lambda_2$.
If $\omega'_1 = \lambda\omega_1$ for some $\lambda \not= 0$,
then $V'_2 = V_2$ as the annihilator of $\omega_1$ or $\omega'_1$.
If $\omega'_1 = \mu\omega_2$ for some $\mu \not= 0$,
then we have $V'_2 = V_1$.
By the symmetric argument, we obtain also that $V'_1 = V_1$ or $V'_1 = V_2$.
Thus we have $V_1' = V_1$, $V_2' = V_2$ or
$V_1' = V_2$,  $V_2' = V_1$.
Assume $V_1' = V_1$, $V_2' = V_2$, then,
restricting $\omega$ to $V = V_1 \oplus V_2$, we have
$\omega'_1 = \omega_1$, $\omega'_2 = \omega_2$.
Assume $V_1' = V_2$,  $V_2' = V_1$,
then we have $\omega'_1 = -\omega_2$, $\omega'_2 = - \omega_1$.
\end{proof}
